\documentclass[a4paper,12pt]{article}

\usepackage[T1]{fontenc}
\usepackage{xltxtra}
\usepackage[francais]{babel}
\usepackage{fancyhdr}

\usepackage{epsfig}
\usepackage{calc}
\usepackage{url}
\usepackage{boxedminipage}

\usepackage{titlesec}
%Pour enlever cette merde de chapter!
%\titleformat{\chapter}[hang]{\bf\huge}{\thechapter}{2pc}{}
\usepackage{graphicx}
\usepackage{color}
\usepackage{float}

\usepackage{tabularx}
\usepackage[hidelinks]{hyperref}
\usepackage[sort, authoryear]{natbib}

% --------------------------------
%Pour compiler sur Ubuntu
%%\usepackage[latin1]{inputenc}   

% --------------------------------
%Pour compiler sur Manjaro        
\usepackage[utf8]{inputenc}
\usepackage{algorithmicx}
\usepackage{algpseudocode}

% --------------------------------
% Fonction of available fonts
%\usepackage{fontspec} 
%\setmainfont{Adobe Garamond Pro}       
%\setmainfont{EB Garamond}
% --------------------------------

%\usepackage[language=french]{csquotes}



%%%%%%%%%%%%%%%%%%%%%%%%%%%%%%%%%%%%%%%%%%%%%%%%%%%%%%%%%%%
%%%%%%%%%%%%%%%%%%%%%%%%%%%%%%%%%%%%%%%%%%%%%%%%%%%%%%%%%%%
%% Définitions à personnaliser 

%% Pour les noms, utilisez la première lettre du prénom suivi du 
%% nom de famille (première lettre majuscule, reste en minuscule).


%%%% Indiquer le nom de l'encadrant ci-dessous:

\def\nomEncad{Martin \textsc{Quinson}}

%% Si le projet est co-encadré indiquer les deux noms à la suite dans 
%% Le même champs


%%%% Indiquer les noms des étudiants participant ci-dessous:

\def\nomEtudA{Louisa \textsc{Bessad}}

%% Si le projet est encadré par moins de 4 étudiants laissez
%% les variables inutiles vides


%%%% Indiquer la référence (numero) et le nom du sujet ci-dessous:

\def\refProjet{Numéro Projet} 
\def\titreProjetCourt{Titre court}
\def\titreProjetLong{Titre long}

%% Le titre court ne doit pas faire plus d'une vingtaine de caractère
%% résumez le à quelques mots essenciels


%%%% Indiquer le type de document et sa version ci-dessous:

\def\typeDoc{Pré-rapport}
 
%% - Rapport intermédaire
%% - Rapport final

%\let\origsec\section
%\renewcommand{\section}[1]{\newpage\origsec{#1}}



%%%%%%%%%%%%%%%%%%%%%%%%%%%%%%%%%%%%%%%%%%%%%%%%%%%%%%%%%%%
%%%%%%%%%%%%%%%%%%%%%%%%%%%%%%%%%%%%%%%%%%%%%%%%%%%%%%%%%%%
%% Définitions à ne pas modifier
 
%%%% ||| Mise en page verticale ||| 
\setlength{\voffset}{-1in} % a4:reste 297mm pour les 5 suivants:
\setlength{\topmargin}{15mm}         % avant l'en-tête
\setlength{\headheight}{20mm}        % hauteur de l'en-tête 
\setlength{\headsep}{10mm}            % entre l'en-tête et le corps
\setlength{\textheight}{220mm}       % hauteur du corps
\setlength{\footskip}{12mm}          % pied de page par rapport au corps 
%\setlength{\footlength}{2em}

%%%%% --- Mise en page horizontale ---
\setlength{\hoffset}{-1in} % a4:reste 210mm 
\setlength{\oddsidemargin}{25mm}     % entre hoffset et le corps
\setlength{\evensidemargin}{25mm}    % entre hoffset et le corps
\setlength{\marginparwidth}{0mm}     % largeur de la marge
\setlength{\marginparsep}{0mm}       % séparateur corps marge
\setlength{\textwidth}{160mm}        % largeur du corps

%\usepackage{fullpage}
%\setlength{\headheight}{20mm}        % hauteur de l'en-tête 
%\setlength{\headsep}{10mm}            % entre l'en-tête et le corps
%\setlength{\textheight}{200mm}
%\setlength{\footskip}{0mm}          % pied de page par rapport au corps 

\def\annee{2015}



%%%%%%%%%%%%%%%%%%%%%%%%%%%%%%%%%%%%%%%%%%%%%%%%%%%%%%%%%%%
%% Début du document

\begin{document}

\selectlanguage{francais}



%%%%%%%%%%%%%%%%%%%%%%%%%%%%%%%%%%%%%%%%%%%%%%%%%%%%%%%%%%%
%% Définition des en-têtes et pied de pages
\pagestyle{fancyplain}
%\fancyhead{}
%\fancyfoot{}
%
\fancyhead[L]{\textsc{Université Pierre et Marie Curie}\\
          Master Informatique\\ Stage M2 \annee \\ \nomEtudA}
\fancyhead[C]{\textbf{Pré-rapport}}%\\\titreProjetCourt}}
\fancyhead[R]{\textsc{LORIA} \\ {\color{white} b} \\ {\color{white} b} \\ \nomEncad}

\fancyfoot[L]{\includegraphics[width=4cm]{Pictures/png/UPMC_sorbonne.png}}
\fancyfoot[C]{\textbf{\thepage/\pageref{fin}}}
\fancyfoot[R]{\includegraphics[width=4cm]{Pictures/loria_logo.jpg}}

%\lhead[\fancyplain{}{\texttt{Université Pierre et Marie Curie}\\
%          Master Informatique\\ UE \textbf{PSAR} fév. \annee \\ \nomEncad}]
%      {\fancyplain{}{\textsc{Université Pierre et Marie Curie}\\
%          Master Informatique\\ UE \textbf{PSAR} \annee \\ \nomEncad}}
%\chead[\fancyplain{}{\textbf{Projet \refProjet\\\titreProjetCourt}}]
%      {\fancyplain{}{\textbf{Projet \refProjet\\\titreProjetCourt}}}
%\rhead[\fancyplain{}{\nomEtudA\\\nomEtudB}]
%      {\fancyplain{}{\nomEtudA\\\nomEtudB}}
%\lfoot[\fancyplain{}{\epsfig{figure=UPMC_sorbonne.eps,width=3cm}}]
%      {\fancyplain{}{\epsfig{figure=UPMC_sorbonne.eps,width=3cm}}}
%\cfoot[\fancyplain{}{\textbf{\thepage/\pageref{fin}}}]
%      {\fancyplain{}{\textbf{\thepage/\pageref{fin}}}}
%\rfoot[\fancyplain{}{\typeDoc}]
%      {\fancyplain{}{\typeDoc}}


%%%%%%%%%%%%%%%%%%%%%%%%%%%%%%%%%%%%%%%%%%%%%%%%%%%%%%%%%%%

      \begin{center}
        \begin{boxedminipage}{12cm}{
            \begin{center}
              ~\\\LARGE\textbf{\titreProjetLong}\\
              ~\\\large Étudiante: \textbf{\nomEtudA,}\\
              ~\\\large Encadrant: \textbf{\nomEncad}\\
              ~
            \end{center}
            }
        \end{boxedminipage}
      \end{center}

%\vfill

\newpage

\tableofcontents
%\vfill
\newpage

%Abstract
\begin{abstract}
  %Intro d'avant
TODO
  %% Dans le cadre de ce stage nous allons nous intéresser aux applications ditribuées à large échelle et comment on peut les tester et évaluer leurs performances via une combinaison d'émulation et de simulation en utilisant SIMGRID et Simterpose qui sont deux projets européens. SIMGRID a été lancé en 1999 pour étudier des algorithmes d'ordonnancement sur des plateformes hétérogènes dans un environnement distribué et faciliter leur programmation. Il fournit les outils de base nécessaire à la simulation de ce type d'applications. Simterpose s'insère dans le projet SIMGRID afin de pouvoir étudier des applications complètes et pas uniquement leur modèle que l'on fournit habituellement en paramètre au simulateur. Le but est de faire de l'émulation en utilisant un simulateur que sera SIMGRID. Puisque nous nous intéressons aux applications distribuées notre émulateur doit pouvoir\textit{i)} exécuter un grand nombre d'instances d'une même application sur un même système afin de pouvoir debugguer, \textit{ii)} évaluer des applications ayant de nombreuses condition d'exécution (simple n\oe ud, réseau complet), \textit{iii)} collecter les informations concernant l'application pendant qu'elle s'exécute.
\end{abstract}
\newpage

% Content
{\color{white} blabla} \vspace{7cm}


Je voudrais remercier mon tuteur Martin Quinson pour m'avoir permis d'effectuer ce stage, durant lequel j'ai beaucoup appris.

Je tiens également à remercier mon encadrant universitaire Sébastien Monnet pour ses précieux conseils.

\newpage
\section{Introduction}

%% intro/objectif: virtualisation légère d'applications distribuées (tester des applications distribuées réelles: test regression et performance, légère car on veut tester des centaines d'instances)

%% Applications: stockage distribué (CEPH, TAHOE/LAFS) et RT event processing (Storm)

Dans le cadre de ce stage, nous allons nous intéresser aux applications
distribuées. Il s'agit d'applications dont une partie ou la totalité des
ressources n'est pas localisée sur la machine où l'application s'exécute, mais
sur plusieurs machines distinctes. Ces dernières communiquent entre elles via le
réseau pour s'échanger les données nécessaires à l'exécution de
l'application. Les applications distribuées ont de nombreux avantages: elles
permettent notamment d'augmenter la disponibilités des données en se les
échangeant, comme les applications Torrent (BitTorrent, $µ$Torrent...). Grâce au
projet BOINC\footnote{\url{https://boinc.berkeley.edu/}} par exemple, on peut
partager la puissance de calcul inutilisée de sa machine. Depuis une dizaine
d'années, la popularité de ces applications distribuées ne cesse de croître. On
peut notamment penser aux applications de stockage de données
LAFS\footnote{\url{https://tahoe-lafs.org/trac/tahoe-lafs}} et
CEPH\footnote{\url{http://ceph.com/}}. La première apporte un stockage robuste
qui préserve les données même si un site est physiquement détruit. La seconde
souhaite fournir performance, fiabilité et scalabilité.  Elles deviennent de
plus en plus complexes avec des contraintes et des exigences de plus en plus
fortes, en particulier au niveau des performances et de l'hétérogénéité des
plateformes et des ressources utilisées. Il est donc de plus en plus
difficile de créer de telles applications (absence d'horloge et mémoire
centrale, deadlock, race, famine) mais aussi de les tester.  En effet, malgré
l'évolution des applications distribuées, les protocoles d'évaluation de leurs
performances n'ont que peu évolués.
\newline

Actuellement, il existe trois façons principales de tester le comportement
d'applications distribuées \citep{gustedt2009experimental}: l'exécution sur
plate-forme réelle, la simulation et l'émulation.

La première solution consiste à exécuter réellement l'application sur un parc de
machines et d'étudier son comportement en conditions réelles. Cela permet de la
tester sur un grand nombre d'environnements. L'outil créé et développé en partie
en France pour cela est \textbf{Grid'5000}\footnote{Infrastructure de 8000 c\oe
  urs répartis dans la France entière crée en
  2005. \\ \url{https://www.grid5000.fr/mediawiki/index.php/Grid5000:Home}}\citep{GRID5000},
un autre outil développé à l'échelle mondiale est
\textbf{PlanetLab} \footnote{Crée en 2002, cette infrastructure de test compte
  aujourd'hui 1340 c\oe urs. \\ \url{http://www.planet-lab.org}}. Néanmoins,
pour mettre en \oe uvre ces solutions complexes, il faut disposer des
infrastructures nécessaires pour effectuer les tests. Il faut également écrire
une application complète capable de gérer toutes ces ressources
disponibles. Enfin, du fait du partage des différentes plateformes entre
plusieurs utilisateurs, les expériences sont difficilement reproductibles.

La seconde solution consiste à faire de la simulation: on modélise ce que l'on
souhaite étudier (application et/ou environnement) via un programme appelé
simulateur. Dans ce cas, pour pouvoir tester des applications distribuées sur un
simulateur, on doit d'abord abstraire l'application ainsi que l'environnement
d'exécution. Pour cela, on identifie les propriétés de l'application et de son
environnement puis on les transforme à l'aide de modèles mathématiques. Ainsi,
on va exécuter dans le simulateur le modèle de l'application et non
l'application réelle, dans un environnement également modélisé. Cette solution
est donc facilement reproductible, plus simple à mette en \oe uvre, et permet de
prédire l'évolution du système étudié grâce à l'utilisation de modèles
mathématiques. De nos jours, les simulateurs tel que
\textbf{SimGrid}\citep{CASANOVA:SimGrid, MARTIN:SimGrid} peuvent simuler des
applications distribuées mettant à contribution des milliers de
n\oe uds. Néanmoins, avec la simulation on ne peut valider qu'un modèle et non
l'application elle-même.

La troisième solution consiste à faire de l'émulation: on exécute réellement
l'application mais dans un environnement virtualisé grâce à un logiciel,
l'émulateur. Ce dernier joue le rôle d'intercepteur pour virtualiser
l'environnement d'exécution.
%On fera ainsi croire à l'application qu'elle s'exécute sur une machine autre que l'hôte.
Cette solution représente un intermédiaire entre la simulation et l'exécution
sur plateforme réelle visant à résoudre les limitations de ces deux
solutions. En effet, les actions de l'application sont réellement exécutées sur
la machine hôte (la machine réelle sur laquelle s'exécute l'émulation) et grâce
à la virtualisation, l'application pense être dans un environnement différent de
la machine réelle. De plus, cela évite d'avoir deux versions de l'application en
terme de code: une pour la simulation et une pour la production. L'émulation
peut-être faite \textit{off-line} (on sauvegarde les actions de l'application
sur disque et on les rejoue plus tard dans un simulateur) ou \textit{on-line}
(on bloque l'application le temps que le temps de réponse de la plateforme
virtualisée soit calculé).

%% Il existe deux types d'émulation pour les applications distribuées; la
%% virtualisation standard et la légère. On parle de virtualisation légère quand on
%% souhaite tester des applications sur une centaine d'instances.

Au cours de ce stage, nous allons nous intéresser plus particulièrement à l'exécution d'applications distribuées arbitraires au dessus de la plateforme de simulation SimGrid. Pour cela, nous allons présenter en section \ref{section:emulation} les méthodes utilisées
pour faire de la virtualisation %% légère: limitation et interception
. Puis en
section \ref{section:sota} nous présenterons les projets permettant de faire de
l'émulation pour tester des applications dans un environnement
distribué. Ensuite, nous expliquerons en section \ref{section:simterpose},
pourquoi dans le cadre du projet Simterpose c'est la virtualisation légère par
interception qui a été choisie et comment elle fonctionne. Pour finir, nous
concluerons en section \ref{section:ccl}.

\newpage
\section{Pourquoi choisir l'émulation simulée}

Il existe trois façons de tester des applications distribuées. La première consiste à exécuter réellement l'application sur un parc de machines et d'étudier le comportement de l'application en temps-réel, ce que fait actuellement \textbf{Grid'5000} {\color{red}mettre citation}. Néanmoins pour mettre en \oe uvre cette solution complexe il faut disposer des architectures nécessaires pour effectuer les tests. De part le partage des différentes plateformes entre divers utilisateurs les expériences ne sont pas forcément reproductibles. Une deuxième solution consiste à simuler l'exécution des applications à l'aide d'un simulateur tel que \textbf{SIMGRID} {\color{red}mettre citation}. On doit alors modéliser l'application ainsi que l'environnements d'exécution grâce à des modèles mathématiques. Le problème avec la simulation est que l'on exécute pas vraiment l'application, on ne peut alors valider qu'un modèle et pas l'application puisqu'on réécrit l'application selon le modèle. Un des buts du projet étant de tester les applicationis sans avoir accès à leur code source, on ne peut donc pas choisir cette solution ou pas toute seule en tout cas. La troisième solution consiste à faire de l'émulation, c'est-à-dire que nous allons exécuter réellement l'applications mais dans en environnement virtuels. Ainsi nous aurons un simulateurs qui gérera l'environnement, l'application qui s'exécutera réellement sur la machine hôte et une API qui fera le lien entre l'application qui s'exécute et l'environnement simulé. On fera ainsi croire à l'application qu'elle s'exécute sur une machine autre que l'hôte. Il existe deux façons de faire de l'émulation: la dégradation et l'interception. Dans la première on rajoute la couche d'émulation au-dessus de la plateforme réelle (comme un hyperviseur pour une VM). Mais cela nous empêche d'émuler des machines plus puissante que l'hôte car le délai de réponse géré par l'émulateur ne peut-être inférieur à celui de l'hôte sinon l'hôte n'a pas le temps de faire les calculs nécessaires à l'application. Cette solution choisie par \textbf{Distem}{\color{red}mettre citation} est donc limitée à la capacité des plateformes à notre disposition. Dans le cas de l'interception pour faire croire à l'application qu'elle s'exécute sur une machine autre que l'hôte on va attraper toutes ses communications via une API du simulateur qui ensuite les transmetra au simulateur. Ce dernier s'occupera de calculer le temps de réponses de ces commumnications en se basant sur les performances de la machine qu'on est en train de simuler. Les calculs seront effectués sur la machine hôte mais le temps de réponse à l'application sera géré par le simulateur en fonction de l'architecture que l'on simule en utilisant le temps d'exécution de la machine hôte injectés dans le simulateur et en comparant les performances des deux architectures. Ainsi le temps de l'application sera celui du simulateur et non le temps réel. Cette solution est implémentée dans \textbf{Simterpose}{\color{red}mettre citation}. Dans notre cas d'émulation simulation, nous allons utiliser SIMGRID comme simulateur et Simterpose comme API de ce simulateur. Nous aurons donc Simterpose qui permettra d'utiliser SIMGRID avec des applications réelles en leur faisant croire qu'elles s'exécutent sur des machines distribuées. Maintenant que nous savons comment nous allons tester les applications distribuées nous allons voir comment fonctionne notre ``émulateur'' Simterpose.

\newpage
\textbf{INTRO TODO}

\newpage

\section{Simterpose: Symboles sur lesquels faire de la médiation}
\label{section:simterpose}
Dans le cadre du projet Simterpose de virtualisation légère et de test
d'applications distribuées, c'est l'émulation par interception qui a été
choisi. En effet le but final étant de pouvoir évaluer n'importe quelle
application distribuée sur n'importe quel type d'architecture, on pourrait se
retrouver à devoir émuler des machines plus puissantes que l'hôte, ce que
l'émulation par dégradation ne permet pas. Pour cela on va utiliser SIMGRID
comme simulateur et Simterpose comme émulateur. Simterpose qui est une API de SIMGRID nous permettra donc d'utiliser le simulateur avec des applications
réelles tout en leur faisant croire qu'elles s'exécutent sur des machines
distinctes. Simterpose étant l'émulateur qui va nous permettre d'intercepter les
communications de l'application avec la machine sur laquelle elle s'exécute et
de faire de la médiation nous allons étudier son fonctionnement et voir quels
outils présentés en section \ref{section:emulation} ont été choisis.

\subsection{Organisation générale}
%schéma tableau

\subsection{Les communications réseaux}
 %-> syscall -> ptrace (full mediation, address translation)

Lorsque ptrace est appelé en entrée ou sortie d'appel système, les modifications
à apporter ne sont pas forcément les mêmes selon qu'il s'agit d'une action
nécessitant l'utilisation du réseau ou non. Dans le cas d'un simple calcul ce
qu'il faut maintenir pour l'application, c'est une vision du temps correspondant
à celle qui s'écoulerait si elle était vraiment sur la machine simulée. Ainsi en
entrée d'appel système on n'a pas besoin de modifier quoique ce soit, par contre
au retour il faut modifier le temps d'exécution du calcul en le remplaçant par
celui calculé par le simulateur.

Dans le cas d'une communication réseau, le but de Simterpose étant de réussir à
simuler un réseau réel sur un réseau local, il faut gérer la transition entre
réseau local et réseau simulé. En effet l'application possède une adresse IP et
des numéros de ports ``virtuels'' qui ne correspondent pas forcément à ceux
attribués dans le réseau local. De plus on ne peut pas se baser uniquement sur
le numéro de \textit{file descriptor} associé à une socket pour identifier deux
entités qui communiquent entre elles. En effet ce \textit{file descriptor} est
unique pour chaque socket d'un processus, mais plusieurs processus peuvent avoir
un même numéro de \textit{file descriptor} pour des sockets de communications
différentes puisque chacune à son propre espace mémoire. Pour pallier à ce
problème on va utiliser en plus du numéro de socket, les adresses IP et les
ports locaux et distants des deux entités qui souhaitent communiquer comme moyen
d'identification. Pour gérer toutes ces modifications deux solutions ont été
proposées lors d'un précédent stage
\cite{interception:GUILLAUME:interception_syscall}: la ``médiation par
traduction d'adresse'' et la ``full médiation''.

{\color{red}schéma}
\paragraph{Traduction d'adresse}
 Avec ce type de médiation on considère que le noyau gère des
 communications. Ainsi en entrée et sortie d'appel système Simterpose va juste
 s'occuper de la transition entre réseau ``virtuel'' et réseau local, en
 utilisant les informations de communications contenues dans la socket. Pour
 cela Simterpose gère un tableau de correspondances, dans lequel pour chaque
 couple adresse IP et des ports ``virtuels'', on a un couple adresse IP et ports
 ``réel'' sur le réseau associé.  De fait en entrée d'un appel système de type
 réseau (bind, connect, accept ...), Simterpose devra remplacer l'adresse et les
 ports ``virtuels'' de l'application par l'adresse et les ports réels sur le
 réseau local, ainsi l'appel système se fera avec une source qui existe
 réellement sur le réseau. Au retour de l'appel système il faudra remodifier les
 paramètres en remettant l'adresse et les ports ``virtuels'' pour que
 l'application pense toujours être dans son environnement simulé.  La limite de
 cette approche est liée au nombre de ports disponibles sur l'hôte.

\paragraph{Full médiation} 
Dans ce cas le noyau ne va plus gérer des communications car nous allons
empêcher l'application de communiquer via des sockets et même d'établir des
connexions avec une autre application.  Puisqu'il n'y a aucune communication, on
n'a pas besoin de gérer de tableau de correspondance d'adresse et de ports et
les applications peuvent conserver les adresses et les ports simulées qu'elles considèrent
comme réels. Quand l'application voudra faire un appel système de type
communication ou connexion vers une autre application, le processus espion de Simterpose qui
sera notifié via ptrace neutralisera l'appel système. Ensuite ce processus en
utilisant ptrace récupérera en lisant dans la mémoire du processus espionné les
informations à envoyer ou récupérer et ira directement lire ou écrire ces
informations dans la mémoire du destinataire.  Même si la ``full médiation''
permet d'éviter les communications réseaux et de conserver des tables de
correspondances, dans le cas d'applications qui communiquent énormément et
utilisent de grosses données elle s'avère moins efficace. En effet les appels à
la mémoire sont bien plus coûteux que les communications réseaux.

\subsection{Les thread}
 %% syscall clone + libcalls

\subsection{Le temps}
 %% -> syscall (- system wide), VDSO-linker (cross process ou VDSO)

\subsection{DNS}
%% libcalls (ne rien rater), config fake (system wide), intercept 53 ( plus dur que nécessaire, port dns autre ou pas)

%\section{Évaluation}
\label{section:evaluation}

intro

\newpage
\section{Conclusion}
Dans ce rapport, nous avons présenté notre objectif ainsi qu'un état de l'art
des différentes approches et outils existants. Pour pouvoir exécuter n'importe
quel type d'applications distribuées dans des conditions environnementales qui
permettront de les étudier, la meilleure solution semble être la virtualisation
par interception. Les projets existants ne permettant pas de résoudre les quatre
problèmes engendrés par cette virtualisation (gestion du temps, des threads, des
communications réseaux et le DNS), un nouveau projet a été lancé.

L'émulateur Simterpose dévelopé au LORIA permet d'exécuter et de tester des
applications distribuées réelles sans disposer de leur code sources dans un
environnement distribué virtuel. Il se base sur la plateforme de simulation
Simgrid pour mettre en place l'environnement d'exécution dans lequel
l'application pensera s'exécuter. Pour maintenir la virtualisation, les actions
des applications sont interceptées et modifiées pour ensuite être exécutées. On utilise SimGid pour calculer la réponse de l'environnement virtuel aux
différentes actions. La solution proposée intercepte les actions à deux niveaux
différents, appels systèmes et des bibliothèques, afin de ne pas en oublier. Simterpose permet également d'injecter diverses fautes dans la
simulation pour avoir une virtualisation plus réaliste.

Néanmoins, il reste à ajouter certaines fonctionalités (DNS, fonctions de temps)
afin de pouvoir vérifier qu'il est possible de mettre en place une
virtualisation par inteception qui apporte une réponse à notre
problématique. Nous projettons par la suite d'évaluer les performances de notre
projet, afin de le modifier pour le rendre plus efficace et de maximiser la
taille des expériences.



\newpage
\label{fin}

% Bib
\bibliographystyle{unsrt}
\bibliography{src/7.Biblio} % The file containing the bibliography

\end{document}
